%! Diagonalizing a Matrix — Unit 6, Lesson 6.2 — Exit Ticket
\documentclass[10pt]{article}
\usepackage{linalg-article}
\usepackage{linalg-boxes}
\newcommand{\TallMath}[1]{$\displaystyle #1\rule[-1.4em]{0pt}{3.2em}$}
\newcommand{\xx}{\mathbf{x}}
\newcommand{\zero}{\mathbf{0}}

\begin{document}

\pageheader{Unit 6, Lesson 6.2}{Exit Ticket}
\namedateperiod

\vspace{0.15in}

No notes. Show your reasoning. Use $A=\begin{bmatrix} 1 & 4 \\ 0 & 3 \end{bmatrix}$ for items 1--2.

\begin{enumerate}[leftmargin=1.8em, itemsep=14pt]
  \item \textbf{Build $X$ and $\Lambda$.} State the two eigenvalues (it is triangular), find an eigenvector
        for each, and write the eigenvector matrix $X$ and the eigenvalue matrix $\Lambda$.
        \vspace{1.7cm}
  \item \textbf{Diagonalize.} Find $X^{-1}$ (Unit~2 formula) and write $A$ as $X\Lambda X^{-1}$.
        \vspace{1.5cm}
  \item \textbf{Explain.} Once you have $A=X\Lambda X^{-1}$, why is $A^{k}=X\Lambda^{k}X^{-1}$ so much easier
        to compute than multiplying $A$ by itself $k$ times? (One or two sentences.)
        \writelines{2}
\end{enumerate}

\end{document}
